\cleardoublepage
\newrefsection
\chapter{文献综述}

\section{背景介绍}

三维视觉的核心任务，是从图像或视频中恢复场景的几何结构、外观信息以及相机位姿，并进一步支撑稳定、真实的新视角合成。随着虚拟现实、增强现实、数字孪生、空间计算与远程协作等应用不断发展，单纯的二维图像记录已经难以满足用户对于沉浸式浏览、自由视点观察和交互式体验的需求。相较于只能从固定机位回看图像或视频，基于三维场景表示的新视角合成能够让用户在虚拟空间中自由改变观察位置，同时尽量保持几何一致性、纹理细节和视角相关外观，因此成为真实场景数字化、数字内容生产和空间媒体构建的重要基础能力。

进一步来看，高保真三维场景资产的价值正在从“面向人观看”扩展到“面向机器使用”。一方面，面向虚拟现实和数字孪生的室内场景模型需要具备可漫游、可渲染、可更新的能力，以服务于空间展示、远程运维、设施管理和沉浸式交互等任务；另一方面，具身智能、机器人学习和自动驾驶等方向也越来越依赖高质量三维资产来构建仿真环境。对于具身智能系统而言，几何准确且外观真实的三维场景有助于开展导航、抓取、交互和任务规划的训练与评测；对于自动驾驶与移动机器人而言，可控且逼真的三维环境能够用于感知算法验证、闭环仿真和数据合成，降低真实采集与反复试验的成本，并缓解仿真与真实之间的域差异。因此，高保真三维重建已经不再只是静态建模问题，而是连接数字孪生、空间智能和智能体训练的重要基础设施。

\section{国内外研究现状}

\subsection{基于稀疏重建的多视图几何方法}

从技术路线的演进来看，早期三维重建长期建立在局部特征匹配与多视图几何之上，其核心目标是从多张图像中恢复相机位姿和稀疏三维结构。典型流程通常先利用 SIFT 进行关键点检测与描述 \cite{lowe2004sift}，再结合特征匹配、几何验证、增量式运动恢复结构（Structure from Motion, SfM）与光束法平差（Bundle Adjustment, BA），对相机位姿和稀疏三维结构进行联合优化 \cite{triggs2000ba,schonberger2016sfm}。这类方法具有几何模型清晰、可解释性强、工程体系成熟、易于与相机标定结合等优点，至今仍然是许多重建系统的基础。

随着深度学习在局部特征领域的发展，学习式前端开始被用于强化这一类稀疏重建管线。以 SuperPoint \cite{detone2018superpoint} 为代表的学习式关键点与描述子方法，在弱纹理和较大视角变化场景中展现出更强的鲁棒性；LightGlue \cite{lindenberger2023lightglue} 则通过自适应匹配机制，在保持较高匹配精度的同时显著提升了稀疏特征匹配效率。将学习式特征与匹配器引入传统 SfM 流程后，可以提升复杂室内环境中的对应关系估计质量，使视频化重建和在线定位在工程上更具可用性。从这个意义上看，深度学习辅助的 SfM 并不是脱离传统多视图几何的新范式，而更像是在稀疏重建框架内部对前端特征提取与匹配能力的递进式增强。

然而，无论是传统手工特征方法，还是引入学习式前端后的 SfM 管线，其本质上都仍属于稀疏重建路线。这一路线的局限性也较为明显。首先，在弱纹理、重复纹理、反光表面、运动模糊以及大畸变成像条件下，对应关系估计仍然可能不稳定，进而影响位姿恢复和稀疏点云质量。其次，这类方法的输出通常是相机参数和稀疏点云，本身并不能直接形成高保真的三维场景表示和新视角合成结果，往往仍需依赖后续的稠密重建、网格生成、纹理映射或辐射场优化步骤。更重要的是，稀疏重建只能提供有限的几何锚点和位姿先验，对于遮挡频繁、结构复杂、纹理分布不均的室内场景而言，若前端恢复结果存在误差，这些误差还会继续传递到后续的显式或隐式表示学习之中。因此，稀疏重建虽然是许多系统的基础，但其自身仍难以单独支撑高保真、强泛化的新视角合成任务，这也正是后续 NeRF、3DGS 及其变体持续发展的重要动因。

\subsection{神经辐射场方法}

为了解决传统几何重建难以直接支撑高质量新视角合成的问题，神经辐射场方法为三维场景表示提供了新的范式。NeRF 将场景表示为连续的隐式辐射场，并通过多层感知机和体渲染过程实现高质量新视角合成 \cite{mildenhall2020nerf}。该类方法能够有效建模视角相关外观和复杂光照效应，在图像保真度方面取得了显著进展，因此迅速成为新视角合成领域的重要方向。

在 NeRF 之后，研究者开始围绕训练和渲染效率展开大量改进。例如，Instant-NGP 利用多分辨率哈希编码显著加速神经场优化 \cite{muller2022instantngp}，Plenoxels 则利用显式稀疏体素与球谐表示替代纯神经网络建模，以降低训练成本 \cite{fridovichkeil2022plenoxels}。这些工作推动了神经场从“高质量”向“更高效率”发展，但从整体上看，基于体渲染的隐式表示在训练时间、推理速度和部署成本上仍然存在一定瓶颈，尤其是在需要实时交互和批量生成三维资产的场景中，其应用仍受到限制。

\subsection{三维高斯泼溅（3DGS）的基础方法}

三维高斯泼溅（3DGS）\cite{kerbl2023} 代表了显式场景表示在新视角合成方向上的重要进展。与 NeRF 通过体渲染积分隐式表示场景不同，3DGS 直接使用一组可学习的三维高斯基元来描述场景。每个高斯通常由中心位置、旋转、尺度、透明度以及颜色属性共同确定，其中颜色通常用球谐系数来建模，以表达一定程度的视角相关外观。若将第 \(i\) 个高斯记为 \(g_i\)，则其核心几何参数可以理解为均值 \(\mu_i\) 与协方差 \(\Sigma_i\)，其中协方差常通过“旋转矩阵 + 对角尺度”进行参数化，从而保证协方差矩阵始终半正定。这种参数化方式使 3DGS 既能表达各向同性的小球形图元，也能表达细长、扁平、具有方向性的局部结构。

在初始化阶段，3DGS 往往依赖 SfM 输出的稀疏点云与相机位姿作为先验。稀疏点为高斯中心提供了初始分布，相机位姿则定义了训练视角下的投影关系。相较于从随机噪声开始优化，这种基于多视图几何的初始化能够显著缩小搜索空间，使后续优化更快进入合理的场景结构区域。也正因为如此，3DGS 与 COLMAP 等传统几何重建工具之间通常具有紧密耦合关系：前者负责恢复较可靠的相机参数和稀疏几何，后者在此基础上进一步优化外观与显式渲染图元。

3DGS 的核心在于其可微渲染机制。对于每个三维高斯，算法首先根据当前相机的投影关系将其映射到像平面上，得到一个二维椭圆高斯；随后按照深度顺序对这些二维高斯进行透明度混合。与传统点云渲染相比，这种椭圆高斯的投影方式天然具有覆盖区域，能够在保持显式表示的同时获得连续的图像合成效果。Kerbl 等人的工作特别强调了可微光栅化的重要性，即利用高斯在屏幕空间中的投影形状近似其对像素的贡献，并通过深度排序和透明度累积实现高质量实时渲染。正是这一设计，使 3DGS 避开了 NeRF 中昂贵的沿光线采样与体积分步骤，在训练和推理阶段都具备更高效率。

3DGS 的优化过程通常围绕“渲染误差最小化 + 结构自适应增删”展开。在损失函数层面，原始 3DGS 主要采用基于重建图像与真值图像之间差异的光度损失，并将 \(L_1\) 误差与结构相似性指标结合起来，从而兼顾局部细节与整体结构一致性 \cite{kerbl2023}。但仅依赖固定数量的高斯往往难以适应真实场景中复杂的几何和纹理分布，因此 3DGS 又引入了增密与剪枝机制：对于投影梯度较大、说明误差敏感的区域，通过克隆或分裂高斯来提升局部表示能力；对于透明度极低、贡献较小或尺度不合理的高斯，则进行裁剪和删除。换言之，3DGS 并不是在一个固定拓扑结构上训练，而是在优化过程中动态调整图元数量和分布。

从算法行为上看，增密机制对 3DGS 的性能至关重要。若高斯数量过少，模型难以表达细节，尤其在纹理边缘、薄结构和遮挡边界附近容易出现模糊与空洞；若高斯数量无控制地快速增加，则会导致冗余图元过多，使模型体积膨胀、训练不稳定并拖慢渲染速度。因此，3DGS 的基础算法可以概括为“显式高斯表示 + 可微光栅化渲染 + 误差驱动的自适应结构优化”。这一组合使它在新视角合成质量、训练速度和实时渲染之间取得了较好的平衡，也为后续几乎所有 3DGS 变体提供了统一的出发点。

\subsection{3DGS 的优化方向}

围绕 3DGS 的后续研究主要聚焦于四类问题：其一是大场景和远近尺度变化下的渲染代价不稳定，需要引入更系统的细节层次（Level of Detail, LoD）机制；其二是原始 3DGS 虽然渲染质量较高，但在几何一致性、抗锯齿和漂浮伪影方面仍存在缺陷，需要进一步提升视觉与几何质量；其三是训练速度仍受优化与光栅化开销制约，需要更高效的训练算法和渲染管线；其四则是原始 3DGS 主要建立在透视相机假设上，面对鱼眼、广角以及其他非线性成像模型时会出现投影误差累积，因此需要把 3DGS 从“透视相机专用表示”推进到“面向通用相机的统一表示”。因此，当前 3DGS 的优化研究已经从单纯提高 PSNR 扩展为对“可扩展性、质量、训练效率和相机建模能力”的系统改进。需要说明的是，由于不同论文所采用的数据集、图像分辨率和统计口径并不完全一致，下文表格主要用于综述性的横向参考，不宜将不同表格中的数值简单视为严格公平的一一对比。

\begin{table}[htbp]
    \centering
    \caption{3DGS 优化方向总览}
    \label{tab:gs-overview}
    \footnotesize
    \setlength{\tabcolsep}{3pt}
    \begin{tabular}{>{\raggedright\arraybackslash}p{2.1cm}>{\raggedright\arraybackslash}p{4.2cm}>{\raggedright\arraybackslash}p{5.2cm}>{\raggedright\arraybackslash}p{2.0cm}}
        \toprule
        优化方向 & 代表方法 & 主要目标 & 常见指标 \\
        \midrule
        LoD 与大场景渲染 & V3DG, Octree-GS, CLoD-GS & 为大场景建立多尺度组织，在远近尺度变化下稳定控制计算量 & PSNR, SSIM, LPIPS, 存储开销 \\
        质量提升 & Mip-Splatting, 2DGS, Geometry-Grounded GS, StableGS & 抑制混叠，增强几何一致性，并减少漂浮伪影等问题 & PSNR, SSIM, LPIPS \\
        训练加速 & 3DGS$^2$, Speedy-Splat & 通过优化器设计或渲染管线改造缩短收敛时间与训练成本 & PSNR, SSIM, LPIPS, 训练时间 \\
        通用相机适配 & 3DGUT, 3DGEER & 使 3DGS 可直接支持鱼眼、广角等非透视成像，并降低边缘视场中的投影误差 & PSNR, SSIM, LPIPS \\
        \bottomrule
    \end{tabular}
\end{table}

\subsubsection{大场景下的细节层次方法}

对于城市级或大范围室内外场景，原始 3DGS 面临一个突出问题：无论当前视点远近如何，位于视锥内的大量高斯都可能参与渲染，导致远景视角下仍然承受过高的计算负担。因此，LoD 成为大场景 3DGS 优化中的关键议题。V3DG 针对这一问题提出了较为典型的后处理式 LoD 思路 \cite{yang2025v3dg}。其核心思想是先离线地将已有 3D 高斯资产组织为不同粒度的簇，并在多个层级上构建可替换的细节版本；在线渲染时，再根据用户给定的屏幕投影面积容忍度选择合适粒度的簇。换言之，V3DG 更接近传统图形学中的资产虚拟化思路：先做多层级简化，再在运行时进行选择。这样做的优势是适合大规模组合场景的实时渲染，但代价是前处理复杂度较高，并且更偏向“离线构建 + 在线选择”的图形系统范式。

与 V3DG 相近但更强调显式层级组织的，是 Octree-GS 提出的具备 LoD 结构的三维高斯表示 \cite{ren2025octreegs}。该方法利用多尺度锚点和渐进式训练将高斯安排到不同 LoD 层级，在远景视角优先使用较粗层级，在近景时再逐步调入更精细的高斯，从而提升大场景中的渲染稳定性与帧率。相比单纯依赖后处理裁剪的做法，这类方法开始把“多尺度高斯组织”直接纳入表示与训练设计之中。

更进一步地，CLoD-GS 则将 LoD 从离散层级推进到连续层级 \cite{cheng2025clodgs}。传统 LoD 常常依赖若干固定级别模型的切换，容易引起细节突变和跳变伪影；而 CLoD-GS 将距离相关的衰减参数直接引入每个高斯，使高斯的不透明度能够随着观察距离连续变化。与此同时，它还设计了虚拟距离缩放和由粗到细的训练策略，使模型在训练阶段就对不同观察尺度具备适应性。这类方法的重要意义在于：LoD 不再只是训练后额外附加的渲染加速手段，而是被内生地整合进 3DGS 的表示与训练过程之中。相关工作的代表性量化结果如表~\ref{tab:gs-lod} 所示。需要指出的是，V3DG 更偏向组合式场景资产虚拟化，其公开结果主要强调簇组织与在线调度，因此这里的表格主要呈现更容易统一到 PSNR、SSIM、LPIPS 与存储开销口径上的 Octree-GS 与 CLoD-GS。总体而言，这一方向主要面向大场景和多尺度浏览任务，对于需要稳定帧率和长期部署的数字孪生系统尤其重要。

\begin{table}[htbp]
    \centering
    \caption{LoD 方法在 Tanks \& Temples 上的代表性结果}
    \label{tab:gs-lod}
    \small
    \begin{tabular}{lcccc}
        \toprule
        方法 & PSNR$\uparrow$ & SSIM$\uparrow$ & LPIPS$\downarrow$ & 存储开销/MB$\downarrow$ \\
        \midrule
        3DGS & 23.70 & 0.853 & 0.169 & 372.19 \\
        Octree-GS & 24.17 & 0.858 & 0.161 & 383.90 \\
        CLoD-GS & 23.75 & 0.853 & 0.170 & 278.53 \\
        \bottomrule
    \end{tabular}
    \caption*{\footnotesize 注：本表统一整理自面向大场景的公开结果，表中的“存储开销”表示论文给出的模型内存或存储规模；V3DG 更侧重在线虚拟化调度，因此未纳入该量化表。}
\end{table}

\subsubsection{3DGS 的质量提升方法}

除了体积和速度问题外，原始 3DGS 在图像与几何质量上也存在若干典型缺陷。首先是混叠问题。原始 3DGS 在不同采样率、不同焦距和不同观察距离下容易出现锯齿、膨胀或细节腐蚀。Mip-Splatting 针对这一问题提出了两项关键改进：一是在三维空间引入平滑滤波器，约束高斯尺寸与输入视角所支持的最高采样频率相匹配；二是在二维图像空间引入多级纹理滤波器，以更合理地替代原始 3DGS 的固定膨胀滤波 \cite{yu2024mipsplatting}。该工作的重要性在于，它将 3DGS 从“单一分辨率下效果较好”推进到“跨尺度、跨焦距更稳定”的无混叠表示。

其次，原始 3DGS 在几何重建方面并不总是理想，因为三维高斯本质上更偏向体元式外观拟合，而不是严格的表面建模。2DGS 因此提出将场景从三维体高斯改为二维定向高斯盘，即把图元直接贴附到局部表面上 \cite{huang2024twodgs}。这种做法的核心收益在于多视图几何更一致，特别适合薄结构、边界和法线敏感区域。为了进一步稳定优化，2DGS 还引入了射线与高斯盘相交、深度畸变约束和法线一致性约束，从而在维持实时渲染能力的同时显著提高几何精度。若只从 PSNR、SSIM 和 LPIPS 观察代表性结果，如表~\ref{tab:gs-quality} 所示，Mip-Splatting 在无混叠新视角合成指标上更突出，而 2DGS 则在保持较强外观质量的同时强调几何一致性。

近期又有工作尝试从更根本的角度为高斯表示提供几何支撑。Geometry-Grounded Gaussian Splatting 通过将高斯基元解释为一种随机实体表示，为高斯场景表示提供了更明确的几何语义基础，并借助其体积性质高效渲染高质量深度图，以支持更鲁棒的形状提取 \cite{zhang2026geometrygroundedgs}。与 2DGS 主要通过改变图元形态来提高几何一致性不同，这类工作试图从理论建模层面回答“高斯到底如何对应显式几何实体”这一问题，因此更偏向于为高质量表面恢复和后续几何处理提供基础。

此外，3DGS 训练中还常见一种称为漂浮伪影的问题，即一些高斯漂浮在非真实表面位置上，却仍然参与图像拟合。StableGS 针对这一问题提出了更稳定的训练框架 \cite{wang2025stablegs}。其核心思想是通过跨视角深度一致性约束限制漂浮伪影的产生，并通过双透明度机制在几何约束与外观细节之间进行解耦，从而在消除漂浮伪影的同时避免严重牺牲外观质量。与前述方法相比，StableGS 的贡献更集中在训练稳定性和几何可信度，而 Geometry-Grounded GS 则更偏向深度图与表面提取质量。因此，这一类工作虽然也会带来一定的 PSNR、SSIM 与 LPIPS 改善，但研究重点已经明显从“单纯提分”转向“几何质量与视觉质量同步优化”。可以看出，当前 3DGS 的“质量提升”已不再局限于提高像素指标，而是开始同时关注无混叠、几何感知和无漂浮伪影三个层面。

\begin{table}[htbp]
    \centering
    \caption{3DGS 质量提升方向的代表性结果}
    \label{tab:gs-quality}
    \small
    \begin{tabular}{llccc}
        \toprule
        方法 & 数据集与设置 & PSNR$\uparrow$ & SSIM$\uparrow$ & LPIPS$\downarrow$ \\
        \midrule
        3DGS & Mip-NeRF360 平均 & 27.21 & 0.815 & 0.214 \\
        Mip-Splatting & Mip-NeRF360 平均 & 27.79 & 0.827 & 0.203 \\
        2DGS & Mip-NeRF360 室内场景 & 30.39 & 0.924 & 0.182 \\
        \bottomrule
    \end{tabular}
    \caption*{\footnotesize 注：该表用于说明质量提升方向的代表性趋势，其中 2DGS 原文更强调几何精度和表面一致性，因此与其他方法并非完全同口径对比。}
\end{table}

\subsubsection{3DGS 的训练加速方法}

虽然 3DGS 相比 NeRF 已经显著降低了训练和渲染开销，但其训练过程仍然往往需要数万次迭代，并依赖增密和频繁光栅化，因此加速问题依然重要。一类思路是改进优化算法本身。3DGS$^2$ 观察到不同高斯属性对图像损失的影响具有较强局部性与弱耦合性，因此可以围绕单个高斯及其属性构造局部的牛顿型近似求解过程，从而得到近似二阶收敛的训练算法 \cite{lan2025dgs2}。与传统基于 SGD 的训练相比，这类方法的价值在于减少达到同等误差所需的迭代次数，使 3DGS 的训练从“分钟级”进一步压缩。

另一类思路则聚焦于渲染与管线效率。Speedy-Splat 指出，原始 3DGS 在渲染阶段存在高斯定位不够精确和图元数量偏多的问题，因此同时从“稀疏像素”和“稀疏图元”两个角度优化整个训练管线 \cite{hanson2025speedysplat}。该方法一方面优化渲染管线对高斯的定位方式，以减少不必要的像素处理；另一方面在训练阶段引入更有效的剪枝策略，从而降低图元规模并同步缩短训练时间。与 3DGS$^2$ 更偏算法优化不同，Speedy-Splat 更偏系统层面的训练-渲染协同加速。表~\ref{tab:gs-train} 给出了这一方向的代表性量化结果。可以看到，3DGS$^2$ 的优势在于以更少迭代达到更高收敛质量，而 Speedy-Splat 则更强调在牺牲有限画质的前提下显著压缩训练时长并提升渲染吞吐。

\begin{table}[htbp]
    \centering
    \caption{3DGS 训练加速方向的代表性结果}
    \label{tab:gs-train}
    \small
    \begin{tabular}{llcccc}
        \toprule
        方法 & 数据集与设置 & PSNR$\uparrow$ & SSIM$\uparrow$ & LPIPS$\downarrow$ & 训练时间 \\
        \midrule
        3DGS 基线 & Mip-NeRF360（3DGS$^2$ 设定） & 29.18 & 0.871 & 0.183 & 1307 秒 \\
        3DGS$^2$ & Mip-NeRF360（3DGS$^2$ 设定） & 29.42 & 0.876 & 0.168 & 5--10$\times$ 更快 \\
        3DGS 基线 & Mip-NeRF360（Speedy-Splat 设定） & 27.55 & 0.814 & 0.222 & 23.2 分钟 \\
        Speedy-Splat & Mip-NeRF360（Speedy-Splat 设定） & 26.94 & 0.782 & 0.296 & 15.7 分钟 \\
        \bottomrule
    \end{tabular}
    \caption*{\footnotesize 注：训练加速论文常在各自设定下报告速度与质量，本表主要用于体现“质量变化与训练时间缩短”的总体趋势。}
\end{table}

\subsubsection{3DGS 的通用相机适配方法}

原始 3DGS 的核心渲染流程建立在透视相机投影之上，即先把三维高斯映射为屏幕空间中的二维椭圆，再进行基于深度排序的 splatting。对于标准针孔相机，这一近似通常较为有效；但在鱼眼、超广角或其他畸变相机中，射线方向与像素坐标之间的关系具有更强的非线性，尤其在图像边缘区域，简单的线性化投影会导致高斯形状、覆盖范围和透明度混合关系出现系统误差。因此，鱼眼适配并不只是“给 3DGS 增加一组新的相机内参”，而是需要重新审视高斯从三维空间到观测平面的投影与积分方式。

3DGUT 是这一方向中较具代表性的工作 \cite{wu2025gut}。该方法的目标，是在尽量保留 3DGS 原有高效 splatting 框架的同时，使其能够支持畸变相机、滚动快门乃至反射、折射等次级射线。其关键思想是将三维高斯通过无迹变换（Unscented Transform）表示为一组 sigma points，再把这些点送入任意给定的相机或射线生成模型中进行传播，以此近似估计非线性投影后的二维分布。这样做的优点在于：相机模型被从“固定透视投影公式”解耦为“任意射线生成器”，因此方法本身具有很强的通用性，也便于接入鱼眼与广角成像模型。与此同时，3DGUT 仍然延续了原始 3DGS 高效的屏幕空间 splatting 路线，因此在工程实现上具有较低的替换成本。不过，由于其本质上仍然依赖 sigma points 对非线性投影进行近似，当视场角很大、边缘畸变很强或高斯尺度较大时，近似误差仍可能累积。

3DGEER 则在此基础上进一步追求“通用相机下的精确渲染” \cite{huang2025geer}。该方法指出，若仍然把通用相机下的高斯贡献压缩为屏幕空间椭圆，就很难彻底消除由非线性投影带来的误差。因此，3DGEER 不再把核心问题表述为“如何近似得到二维椭圆”，而是直接在任意相机射线下计算三维高斯沿光线的体积分贡献，从而使渲染过程与具体相机模型解耦。为了避免精确积分带来的计算代价过高，该方法又进一步设计了 Particle Bounding Frustum 和 Bipolar Equiangular Associative Projection（BEAP）等加速结构，用于更高效地完成高斯与通用射线之间的关联。与 3DGUT 相比，3DGEER 的研究重点已经从“在原有 splatting 框架中兼容非透视相机”转向“在通用射线空间中重新定义高斯渲染算子”，因此更加适合分析鱼眼和大视场条件下边缘区域的重建误差来源。

从研究脉络上看，3DGUT 和 3DGEER 代表了鱼眼适配的两条不同技术路线：前者强调在保留 3DGS 原始效率优势的前提下对非线性投影做高效近似，后者则更进一步，直接从渲染积分层面消除投影近似误差。表~\ref{tab:gs-fisheye} 给出了在 ScanNet++ 全视场设定下的代表性结果。可以看到，在统一评测口径下，3DGUT 已经将 PSNR 提升到 30.64，但 3DGEER 进一步把 PSNR、SSIM 和 LPIPS 分别提升到 31.50、0.953 和 0.126，说明当任务真正转向鱼眼和广角相机时，投影建模精度本身会直接影响最终重建质量。对本课题而言，这一方向的启发在于：若希望尽量复用现有透视 3DGS 管线，3DGUT 提供了一条较低改造成本的通用相机兼容路径；而若目标是在鱼眼全视场尤其是边缘区域追求更高保真度，则 3DGEER 所体现的“精确射线积分”思想更值得重点关注。

\begin{table}[htbp]
    \centering
    \caption{通用相机适配方法在 ScanNet++ 全视场上的代表性结果}
    \label{tab:gs-fisheye}
    \small
    \begin{tabular}{lccc}
        \toprule
        方法 & PSNR$\uparrow$ & SSIM$\uparrow$ & LPIPS$\downarrow$ \\
        \midrule
        FisheyeGS & 27.81 & 0.946 & 0.139 \\
        EVER & 29.47 & 0.925 & 0.167 \\
        3DGUT & 30.64 & 0.944 & 0.150 \\
        3DGEER & 31.50 & 0.953 & 0.126 \\
        \bottomrule
    \end{tabular}
    \caption*{\footnotesize 注：本表采用 3DGEER 补充材料中对 ScanNet++ 五个场景的全视场统一评测结果，用于保证不同方法之间的指标口径一致。}
\end{table}

\subsection{基于 Transformer 的直接高斯生成方法}

除“对单个场景逐场优化”这一主流范式外，近年来还出现了另一条重要路线，即利用大规模数据训练通用模型，使其能够从少量输入图像中直接预测高斯场景表示。这类方法的基本目标，是把 3DGS 从“每个场景单独训练”推进到“面向未知场景的前馈式重建器”。GS-LRM 是其中较早的代表之一 \cite{zhang2024gslrm}。该方法基于 Transformer 架构，将多视图图像切分为图像块后送入统一网络，并直接为每个像素预测对应的高斯参数，再将这些高斯融合为完整场景。值得注意的是，GS-LRM 在对象级任务中使用了 Objaverse 等大规模三维数据，在场景级任务中则使用了 RealEstate10K 等真实视频数据，这说明“大数据预训练 + 前馈式高斯预测”已经成为一个独立研究方向。不过，GS-LRM 仍然依赖输入位姿，因此更接近“高效大模型重建器”，而非真正的无位姿方法。

在无位姿输入的条件下，Splatt3R 代表了更加直接的无位姿高斯生成路线 \cite{smart2024splatt3r}。该方法建立在 MASt3R 这类基础几何重建模型之上，在保留其跨图特征交互能力的同时，额外预测构成高斯图元所需的协方差、球谐系数、不透明度和偏移量等属性，从而直接从两张未知相机参数的图像构造三维高斯泼溅图元。Splatt3R 的一个重要设计是“先几何、后渲染”的训练思路，即先利用三维点几何损失稳定空间结构，再加入新视角合成损失优化外观，这在一定程度上缓解了直接从无位姿图像训练 GS 时容易陷入局部最优的问题。

NoPoSplat 则进一步将无位姿 3DGS 推向更强的通用化 \cite{ye2024noposplat}。其核心思想是以一个输入视图的局部坐标系作为规范空间，并让网络直接在该空间中预测所有输入视图对应的高斯表示，从而规避了先估计全局位姿、再做坐标变换所带来的误差传播。同时，该方法引入相机内参标记，将尺度信息以标记形式并入 Transformer 输入，缓解了无位姿情况下的尺度歧义问题。实验结果表明，这种规范空间直接建模方式在低重叠视图下尤其有效。

SelfSplat 则强调在无位姿、无显式三维先验条件下的泛化重建 \cite{kang2024selfsplat}。它将自监督深度估计、匹配感知的位姿网络和高斯属性预测统一到同一框架中，使相机位姿估计与显式高斯重建能够相互促进。与 Splatt3R 更依赖已有几何基础模型不同，SelfSplat 试图在端到端框架中完成深度、位姿和高斯表示的联合推断，因此在方法论上更接近“从输入图像中自举出几何与外观”的可泛化 3DGS。

总体而言，这一方向说明 3DGS 已经不再仅仅是一种场景级优化表示，而正在向“可泛化、可前馈、可无位姿输入”的重建模型演化。对本课题而言，这类方法的启发在于：未来鱼眼室内重建并不一定只能依赖严格的逐场景优化流程，也可以借助大规模数据预训练和 Transformer 式跨视图建模，使系统在未知环境、弱位姿先验乃至快速部署场景中具备更强的适应能力。当然，这一类方法当前在超高保真重建、极端畸变相机适配以及室内细节恢复方面仍存在不足，因此更适合作为长期发展方向，而非直接替代当前基于标定和显式优化的高保真重建流程。

\subsection{当前研究的局限性}

尽管上述方法在新视角合成质量、训练效率和大场景扩展性方面已经取得了显著进展，但从“真实室内高保真重建”的需求出发，现有研究仍存在较为明显的局限。首先，对于以 NeRF、3DGS 及其变体为代表的基于优化的方法而言，其训练目标在很大程度上仍然由 \(L_1\) 损失、SSIM 以及类似的光度损失主导 \cite{mildenhall2020nerf,kerbl2023}。这类损失的优势在于能够直接约束渲染图像与观测图像之间的像素一致性，但它们主要关注的是“已观测视图是否拟合得足够好”，而不是“场景几何是否被正确恢复”。即便 2DGS、StableGS 和 Geometry-Grounded GS 等工作已经尝试通过表面建模、深度一致性和几何解释增强约束 \cite{huang2024twodgs,wang2025stablegs,zhang2026geometrygroundedgs}，整体训练范式仍然没有彻底摆脱对光度监督的依赖。

这种局限在室内场景中表现得尤为明显。室内环境往往存在遮挡频繁、薄结构丰富、重复纹理较多以及大面积弱纹理墙面等特点，若几何约束不足，模型就更容易通过“在观测视图上看起来正确”的方式拟合外观，而不是恢复真实稳定的三维结构。其直接后果是：对于训练图像分布之外的外插视角，重建结果容易出现结构漂移、局部变形、纹理拉伸和几何不连续等问题。换言之，优化式方法通常在插值视角上表现较好，但在外插视角上的泛化质量较差。对真实室内扫描任务而言，这意味着若想获得较为稳定的重建效果，往往需要更高密度的图像采样和更精细的采集轨迹设计，需要从多个高度、多个朝向和多个距离重复覆盖关键区域，采集成本和实施复杂度都会明显上升。

另一方面，前馈式生成网络虽然试图绕开逐场景长时间优化的问题，但其能力高度依赖大规模训练数据及其分布覆盖范围 \cite{zhang2024gslrm,smart2024splatt3r,ye2024noposplat,kang2024selfsplat}。当前常见训练数据虽然已经涵盖 Objaverse、RealEstate10K 等对象级或视频级资源，但与真实室内场景之间仍然存在显著域差异，例如拍摄设备不同、镜头畸变不同、纹理分布更复杂、光照变化更强以及杂乱遮挡更多。特别是对于本课题所涉及的鱼眼相机输入，大视场角和强畸变会进一步放大这种训练测试分布不一致的问题。因此，前馈式方法在真实场景上的生成质量往往仍不稳定，容易出现重影、双边缘、局部错位和细节模糊等现象；当输入视图重叠不足、位姿未知或存在较强动态物体干扰时，这些问题会更加突出。

总体来看，当前研究的核心矛盾可以概括为：优化式方法具有较好的单场景拟合能力，但为了弥补几何约束不足，需要较密集且设计良好的采集数据；前馈式方法具有更强的推理效率和部署潜力，但受限于训练数据规模与域泛化能力，在真实高保真重建任务中仍难以稳定替代逐场景优化流程。这一局限也恰恰说明，面向鱼眼室内场景的高保真重建，未来仍需要在“几何约束更强的优化框架”和“具有更好泛化能力的前馈模型”之间寻找新的平衡点。

\section{研究展望}

综合现有研究可以看出，三维重建的发展正在从“依赖稠密观测的逐场景优化”逐步走向“引入更强先验的高效生成与重建”。围绕 3DGS 这一当前主流显式表示，未来一个重要方向是缓解其对大量多视角 RGB 图像的刚性依赖。现有优化式 3DGS 往往只有在输入视角覆盖充分时才能获得较稳定的新视角质量，这使得高保真重建仍然高度依赖繁琐的数据采集。未来研究很可能会更多地借助视频生成模型或视图生成模型，通过学习到的时空先验和几何先验，在少量图像输入的基础上生成具有跨视图一致性的补充观测，从而为 3DGS 提供更丰富的训练约束。若这一方向能够较好地解决“生成视图的几何一致性”和“生成误差对后续优化的影响”两个问题，那么稀疏输入下的高质量 3DGS 重建将具备更强可行性。

另一个重要趋势，是提升 3DGS 及类似显式表示对相机位姿噪声的鲁棒性。现有方法通常依赖较高精度的相机外参与内参，位姿一旦存在偏差，3DGS 强大的外观拟合能力便可能把这种误差吸收到场景表示之中，进而在非训练视角下表现为重影、双边缘、漂浮伪影以及局部结构漂移等问题。因此，未来研究不仅需要提高前端位姿估计的精度，还需要在表示与优化阶段引入更有效的鲁棒约束，例如联合位姿优化、跨视角几何一致性约束、深度或法线先验、遮挡感知建模以及针对伪影的正则化设计。换言之，三维重建的后续发展不应只追求更强的拟合能力，还应同时强调“在位姿存在误差时仍然保持结构稳定”的能力。

此外，前馈式高斯生成模型的发展仍然高度受限于高质量三维数据的规模和多样性。Splatt3R、NoPoSplat、SelfSplat 等方法已经展示了从少量图像直接生成高斯表示的潜力，但其泛化能力仍在很大程度上取决于训练数据的覆盖范围。未来如果能够持续构建更高质量、更大规模、更具多样性的三维数据集，包括更准确的多视图图像、相机参数、深度信息、表面几何以及显式高斯表示标注，那么前馈式模型就有望学习到更加稳健的几何先验和外观先验，并显著提升在真实复杂场景中的重建质量。由此可见，三维重建领域的进一步突破，不仅依赖新的网络结构或新的表示方式，也同样依赖系统性的数据建设与更严格的评测基准。

\newpage
\begingroup
    \linespreadsingle{}
    \printbibliography[title={参考文献}]
\endgroup
